\documentclass[12pt,a4paper]{report}

% =========================================================
% PACKAGES
% =========================================================

\usepackage[a4paper,margin=0.85in]{geometry}
\usepackage{graphicx}
\usepackage{float}
\usepackage{caption}
\usepackage{titlesec}
\usepackage{fancyhdr}
\usepackage{hyperref}
\usepackage{enumitem}
\usepackage{setspace}
\usepackage{needspace}

% =========================================================
% GENERAL FORMATTING
% =========================================================

\setstretch{1.08}

\setlength{\parindent}{0pt}
\setlength{\parskip}{3pt}

% =========================================================
% HYPERLINKS
% =========================================================

\hypersetup{
    colorlinks=true,
    linkcolor=black,
    urlcolor=blue,
    citecolor=black
}

% =========================================================
% CHAPTER FORMAT
% =========================================================

\titleformat{\chapter}
  {\normalfont\Large\bfseries}
  {\thechapter}
  {0.8em}
  {}

\titlespacing*{\chapter}
  {0pt}
  {0pt}
  {10pt}

% =========================================================
% SECTION FORMAT
% =========================================================

\titleformat{\section}
  {\normalfont\large\bfseries}
  {\thesection}
  {0.7em}
  {}

\titlespacing*{\section}
  {0pt}
  {7pt}
  {3pt}

% =========================================================
% SUBSECTION FORMAT
% =========================================================

\titleformat{\subsection}
  {\normalfont\normalsize\bfseries}
  {\thesubsection}
  {0.7em}
  {}

\titlespacing*{\subsection}
  {0pt}
  {5pt}
  {2pt}

% =========================================================
% FIGURE CAPTIONS
% =========================================================

\captionsetup{
    font=small,
    labelfont=bf,
    justification=centering,
    skip=3pt
}

% =========================================================
% LIST FORMATTING
% =========================================================

\setlist[itemize]{
    leftmargin=1.4em,
    itemsep=1pt,
    topsep=2pt
}

\setlist[enumerate]{
    leftmargin=1.4em,
    itemsep=1pt,
    topsep=2pt
}

% =========================================================
% HEADER / FOOTER
% =========================================================

\pagestyle{fancy}

\fancyhf{}

\fancyhead[L]{MeterIQ}
\fancyhead[R]{Prototype Report}

\fancyfoot[C]{\thepage}

\renewcommand{\headrulewidth}{0.4pt}

% =========================================================
% DOCUMENT
% =========================================================

\begin{document}


% =========================================================
% TITLE PAGE
% =========================================================

\begin{titlepage}

\centering

\vspace*{1.5cm}

{\Large\bfseries UCS503 -- SOFTWARE ENGINEERING\par}

\vspace{1cm}

{\Huge\bfseries MeterIQ\par}

\vspace{0.3cm}

{\Large Explainable Smart Meter Analytics Platform\par}

\vspace{1.2cm}

{\large\bfseries Prototype Report\par}

\vspace{1.2cm}

{\large
\textbf{Yashit Arora}\\
Roll No.: 1024160006
\par}

\vspace{0.5cm}

{\large
\textbf{Aarushi Gahlawat}\\
Roll No.: 1024160008
\par}

\vfill

{\large
Department of Computer Science and Engineering\\
Thapar Institute of Engineering and Technology\\
September 2026
\par}

\end{titlepage}


% =========================================================
% CHAPTER 1 -- INTRODUCTION
% =========================================================

\chapter{Introduction}

% ---------------------------------------------------------
% BACKGROUND
% ---------------------------------------------------------

\section{Background}

Electricity distribution networks generate large amounts of smart meter data
at regular time intervals. Analysing this data can help identify unusual
consumption behaviour and support better understanding of electricity usage.

MeterIQ is an explainable smart meter analytics platform designed to analyse
smart meter consumption data. The system focuses on two major analytical
functionalities: anomaly or electricity-theft-risk detection and short-term
load forecasting.

The platform is being developed using smart meter data from Mathura and
Bareilly. Instead of directly declaring a meter as a theft case, MeterIQ
identifies unusual patterns and provides information that can support further
human investigation.

% ---------------------------------------------------------
% PROBLEM STATEMENT
% ---------------------------------------------------------

\section{Problem Statement}

Traditional analysis of smart meter data can make it difficult to identify
unusual consumption patterns efficiently. Large datasets also make manual
inspection time-consuming.

The proposed system aims to provide an integrated platform where smart meter
data can be processed, analysed and visualised through a dashboard. The
system also aims to provide explainable information along with detected
anomalies so that the results can be understood by users.

% ---------------------------------------------------------
% OBJECTIVES
% ---------------------------------------------------------

\section{Objectives}

The main objectives of MeterIQ are:

\begin{itemize}
    \item Analyse smart meter consumption data.
    \item Detect unusual consumption patterns and generate anomaly or
          theft-risk indicators.
    \item Forecast short-term electricity load.
    \item Provide explanations for detected unusual behaviour.
    \item Present analytical results through an interactive dashboard.
    \item Support human investigation rather than automatically accusing a
          consumer of electricity theft.
\end{itemize}


% =========================================================
% CHAPTER 2 -- SYSTEM DESIGN
% =========================================================

\chapter{System Design}

% ---------------------------------------------------------
% SYSTEM OVERVIEW
% ---------------------------------------------------------

\section{System Overview}

The MeterIQ system follows a data-driven workflow. Smart meter data is first
prepared and transformed into suitable analytical features. The prepared
data is then used by the anomaly detection and forecasting components. The
resulting information is presented through the dashboard.

The overall workflow is:

\begin{center}
\textbf{Smart Meter Data}
$\rightarrow$
\textbf{Data Preparation}
$\rightarrow$
\textbf{Feature Engineering}
$\rightarrow$
\textbf{Analytics}
$\rightarrow$
\textbf{Dashboard}
\end{center}

The analytics layer consists of two primary modules:

\begin{enumerate}
    \item Anomaly / electricity-theft-risk detection.
    \item Short-term load forecasting.
\end{enumerate}

% ---------------------------------------------------------
% USE CASE DIAGRAM
% ---------------------------------------------------------

\needspace{0.55\textheight}

\section{Use Case Diagram}

The Use Case Diagram describes the major interactions between the
Smart-Meter Data Source, User and the MeterIQ system.

\begin{figure}[H]
    \centering
    \includegraphics[
        width=0.68\textwidth,
        height=0.34\textheight,
        keepaspectratio
    ]{Use Case Diagram.draw.io.pdf}
    \caption{MeterIQ Use Case Diagram}
    \label{fig:usecase}
\end{figure}


% =========================================================
% CHAPTER 3 -- SYSTEM MODELLING
% =========================================================

\chapter{System Modelling}

% ---------------------------------------------------------
% DFD LEVEL 0
% ---------------------------------------------------------

\needspace{0.42\textheight}

\section{Data Flow Diagram -- Level 0}

The Level 0 Data Flow Diagram provides a high-level representation of the
MeterIQ system and shows the interaction between the system and its external
entities.

\begin{figure}[H]
    \centering
    \includegraphics[
        width=0.70\textwidth,
        height=0.18\textheight,
        keepaspectratio
    ]{DataFlowDiagram_level0.pdf}
    \caption{MeterIQ DFD Level 0}
    \label{fig:dfd0}
\end{figure}

% ---------------------------------------------------------
% DFD LEVEL 1
% ---------------------------------------------------------

\needspace{0.47\textheight}

\section{Data Flow Diagram -- Level 1}

The Level 1 DFD provides a detailed representation of the major processes
involved in smart meter data analysis.

\begin{figure}[H]
    \centering
    \includegraphics[
        width=0.78\textwidth,
        height=0.28\textheight,
        keepaspectratio
    ]{DataFlowDiagramLevel1.drawio.pdf}
    \caption{MeterIQ DFD Level 1}
    \label{fig:dfd1}
\end{figure}

% ---------------------------------------------------------
% DFD LEVEL 2
% ---------------------------------------------------------

\needspace{0.47\textheight}

\section{Data Flow Diagram -- Level 2}

The Level 2 DFD further decomposes the anomaly and theft-risk detection
process and represents the flow of information between its components.

\begin{figure}[H]
    \centering
    \includegraphics[
        width=0.78\textwidth,
        height=0.28\textheight,
        keepaspectratio
    ]{DataFlowDiagLevel2.drawio.pdf}
    \caption{MeterIQ DFD Level 2}
    \label{fig:dfd2}
\end{figure}

% ---------------------------------------------------------
% ER DIAGRAM
% ---------------------------------------------------------

\needspace{0.52\textheight}

\section{Entity Relationship Diagram}

The Entity Relationship Diagram represents the major data entities and their
relationships used in the MeterIQ system.

\begin{figure}[H]
    \centering
    \includegraphics[
        width=0.72\textwidth,
        height=0.35\textheight,
        keepaspectratio
    ]{ER_Diagram_MeterIQ.pdf}
    \caption{MeterIQ Entity Relationship Diagram}
    \label{fig:erd}
\end{figure}

% ---------------------------------------------------------
% ACTIVITY DIAGRAM
% ---------------------------------------------------------

\needspace{0.55\textheight}

\section{Activity Diagram}

The Activity Diagram represents the sequence of activities performed during
the MeterIQ analysis workflow, from receiving meter data to generating
analytical results.

\begin{figure}[H]
    \centering
    \includegraphics[
        width=0.72\textwidth,
        height=0.38\textheight,
        keepaspectratio
    ]{Activity_Diagram_MeterIQ.pdf}
    \caption{MeterIQ Activity Diagram}
    \label{fig:activity}
\end{figure}


% =========================================================
% CHAPTER 4 -- PROTOTYPE IMPLEMENTATION
% =========================================================

\chapter{Prototype Implementation}

% ---------------------------------------------------------
% DATA PREPARATION
% ---------------------------------------------------------

\section{Data Preparation}

The smart meter data is prepared before performing analytical operations.
The preprocessing stage includes handling the available meter readings,
organising the time information and preparing the data for further analysis.

Feature engineering is then performed to generate useful temporal and
consumption-related features. These features are used by the analytical
modules.

% ---------------------------------------------------------
% ANOMALY DETECTION
% ---------------------------------------------------------

\section{Anomaly and Risk Detection}

The anomaly detection component is intended to identify unusual patterns in
smart meter consumption data.

The system does not directly classify a consumer as committing electricity
theft. Instead, it produces an anomaly or risk indicator based on unusual
behaviour. The result can then be examined by an authorised human user.

The explainability component is intended to provide information about why a
reading or pattern has been identified as unusual.

% ---------------------------------------------------------
% FORECASTING
% ---------------------------------------------------------

\section{Short-Term Load Forecasting}

The forecasting component is designed to estimate electricity load for a
short-term horizon. A same-hour-last-week baseline is used as an initial
reference for evaluating forecasting performance.

The forecasting workflow uses chronological data separation so that future
information is not used during the training process.

% ---------------------------------------------------------
% DASHBOARD
% ---------------------------------------------------------

\section{Dashboard}

A web-based dashboard is being developed to provide a user-friendly
interface for MeterIQ. The dashboard is intended to allow users to search and
inspect meters, view consumption patterns, examine anomaly or risk
information and view load forecasts.

The dashboard provides a single interface for presenting the results of the
analytical pipeline.


% =========================================================
% FRONTEND PROTOTYPE
% =========================================================

\needspace{0.25\textheight}

\section{Frontend Prototype}

The MeterIQ frontend provides an interactive interface for smart meter
monitoring, analysis and decision support. The prototype includes a login
interface, an operations dashboard and a short-term load forecasting
interface.

% ---------------------------------------------------------
% LOGIN INTERFACE
% ---------------------------------------------------------

\needspace{0.45\textheight}

\subsection{Login Interface}

The login interface provides access to the MeterIQ platform and supports
different prototype roles, including Supervisor and Data Analyst.

\begin{figure}[H]
    \centering
    \includegraphics[
        width=0.50\textwidth,
        height=0.34\textheight,
        keepaspectratio
    ]{login.png}
    \caption{MeterIQ Login Interface}
    \label{fig:login}
\end{figure}

% ---------------------------------------------------------
% OPERATIONS DASHBOARD
% ---------------------------------------------------------

\needspace{0.52\textheight}

\subsection{Operations Dashboard}

The Operations Dashboard provides an overview of smart meter activity. It
presents total smart meters, total consumption, detected anomalies,
high-risk meters and active regions. It also provides consumption trends and
anomaly-risk distribution for operational decision support.

\begin{figure}[H]
    \centering
    \includegraphics[
        width=0.86\textwidth,
        height=0.39\textheight,
        keepaspectratio
    ]{dashboard.png}
    \caption{MeterIQ Operations Dashboard}
    \label{fig:dashboard}
\end{figure}

% ---------------------------------------------------------
% FORECASTING INTERFACE
% ---------------------------------------------------------

\needspace{0.52\textheight}

\subsection{Short-Term Load Forecasting Interface}

The forecasting interface allows a smart meter and machine learning model to
be selected along with the forecast horizon. It presents forecasting
evaluation metrics and the predicted load profile for the selected meter.

\begin{figure}[H]
    \centering
    \includegraphics[
        width=0.86\textwidth,
        height=0.39\textheight,
        keepaspectratio
    ]{forecasting.png}
    \caption{MeterIQ Short-Term Load Forecasting Interface}
    \label{fig:forecasting}
\end{figure}


% =========================================================
% CHAPTER 5 -- CURRENT PROGRESS
% =========================================================

\chapter{Current Progress}

The current prototype development includes the following completed or
developed components:

\begin{itemize}
    \item Problem definition and system requirements.
    \item Use Case Diagram.
    \item DFD Level 0, Level 1 and Level 2.
    \item Entity Relationship Diagram.
    \item Activity Diagram.
    \item Exploratory Data Analysis.
    \item Data preprocessing.
    \item Feature engineering.
    \item Initial dashboard development.
    \item System workflow and prototype design.
\end{itemize}

% ---------------------------------------------------------
% PROTOTYPE STATUS
% ---------------------------------------------------------

\section{Prototype Status}

The prototype currently demonstrates the planned system structure and
dashboard workflow. The final analytical results and model evaluation will
be added after completion and integration of the machine learning modules.

The current report represents the mid-semester prototype stage of MeterIQ
rather than the final project evaluation.


% =========================================================
% CHAPTER 6 -- FUTURE WORK
% =========================================================

\chapter{Future Work}

The planned future work for MeterIQ includes:

\begin{itemize}
    \item Complete the anomaly and electricity-theft-risk detection module.
    \item Complete and evaluate the short-term load forecasting module.
    \item Integrate the analytical models with the dashboard.
    \item Add explainable anomaly information to the user interface.
    \item Perform complete system testing.
    \item Evaluate the system using appropriate performance metrics.
    \item Improve the dashboard based on testing and user feedback.
\end{itemize}


% =========================================================
% BIBLIOGRAPHY
% =========================================================

\begin{thebibliography}{9}

\bibitem{smartmeter}
Smart meter consumption data used for the MeterIQ project,
including data from Mathura and Bareilly.

\bibitem{softwareengineering}
UCS503 Software Engineering course material,
Thapar Institute of Engineering and Technology.

\bibitem{github}
MeterIQ project repository:
\url{https://github.com/aarushi2006/ucs503p-202627odd-MeterIQ}

\end{thebibliography}


% =========================================================
% END DOCUMENT
% =========================================================

\end{document}